
O dimensionamento de fundações rasas constitui uma etapa fundamental no projeto de estruturas de concreto armado, sendo responsável por assegurar a transferência adequada das ações da superestrutura ao maciço de solo de apoio. O desempenho dessas fundações está diretamente associado à segurança global da edificação, à limitação de recalques diferenciais e à preservação da integridade de construções vizinhas \cite{SantosCorrêa}. No contexto brasileiro, o avanço do processo de urbanização tem imposto desafios adicionais ao projeto de fundações, em especial em áreas consolidadas, caracterizadas por terrenos reduzidos, geometrias irregulares e edificações implantadas em proximidade imediata. Dados de 2022 indicam que 87,4 \% da população brasileira residia em áreas urbanas, representando um acréscimo de 16,6 milhões de habitantes em relação a 2010 \cite{IBGE2022_Censo}, o que intensifica a verticalização das construções e a complexidade das soluções de fundação adotadas.

Entre as alternativas de fundações rasas, as sapatas isoladas apresentam ampla aplicação em edificações de pequeno e médio porte, em função de sua simplicidade construtiva e viabilidade econômica. Entretanto, o dimensionamento dessas fundações, na prática corrente, ainda é predominantemente conduzido por procedimentos empíricos e iterativos, baseados em tentativas sucessivas até o cumprimento dos critérios normativos. Tal abordagem depende fortemente da experiência do projetista e, em geral, conduz a soluções conservadoras, com consumo elevado de concreto e consequente aumento dos custos de execução \cite{Moraesetal}. Além disso, os métodos tradicionais raramente integram, de forma sistemática e simultânea, as restrições estruturais, geotécnicas e geométricas, o que se configura como uma limitação crítica em projetos urbanos com múltiplas interferências espaciais entre elementos de fundação.

Além das dificuldades técnicas, o setor da construção civil enfrenta pressões crescentes para reduzir sua pegada de carbono. O concreto é um dos materiais de maior impacto ambiental, sendo que apenas a descarbonatação do calcário durante a produção do cimento responde por cerca de 6 \% das emissões nacionais de dióxido de carbono \cite{SantoroKripka}. O transporte de materiais e as etapas construtivas também contribuem significativamente para esse cenário. Assim, o desenvolvimento de projetos de fundações que utilizem menor volume de concreto representa uma estratégia direta de mitigação dos impactos ambientais, sem comprometer os requisitos de segurança e desempenho estrutural.

Diante desse contexto, ferramentas computacionais de otimização têm se mostrado alternativas promissoras para o dimensionamento de elementos estruturais. Essas ferramentas permitem automatizar análises complexas e integrar critérios de projeto em um único modelo. Estudos recentes \cite{Mora_et_al, Moraesetal, TRENTINI} evidenciam que o uso de técnicas de otimização conduz a soluções mais eficientes do que as obtidas por métodos convencionais, possibilitando reduções expressivas no consumo de concreto e nas emissões de carbono, além de ampliar a vida útil das estruturas. Essas metodologias também favorecem a padronização do processo de dimensionamento e reduzem a subjetividade associada à experiência individual do projetista.

Entre as técnicas de otimização disponíveis, as metaheurísticas destacam-se pela capacidade de explorar amplos espaços de busca e lidar com problemas altamente não lineares. O Algoritmo Genético (AG), em particular, tem demonstrado excelente desempenho em aplicações de engenharia estrutural, apresentando alta precisão e convergência eficiente com esforço computacional reduzido \cite{DosSantos2013}. Essa abordagem permite avaliar simultaneamente múltiplas variáveis e restrições, conciliando critérios de segurança, economia e sustentabilidade. Além disso, o uso de algoritmos evolutivos favorece a incorporação de diferentes cenários de carregamento e incertezas associadas às propriedades do solo e dos materiais, ampliando a robustez das soluções obtidas.

Assim, a aplicação de estratégias de otimização baseadas em metaheurísticas representa um avanço relevante para o dimensionamento geométrico de sapatas isoladas. A integração entre parâmetros estruturais, geotécnicos e normativos em um ambiente computacional possibilita identificar dimensões economicamente ótimas, reduzindo o consumo de materiais e os impactos ambientais, sem comprometer a segurança da estrutura.

Dessa forma, o presente trabalho propõe uma metodologia de otimização mono-objetivo, fundamentada em metaheurísticas evolutivas, para o dimensionamento automático de sapatas isoladas rígidas. O modelo desenvolvido integra, de forma  simultânea,  restrições geométricas, geotécnicas e de resistência previstas nas normas técnicas brasileiras. O objetivo consiste em minimizar o volume total de concreto empregado e, consequentemente, os custos e os impactos ambientais associados. Com isso, pretende-se contribuir para o avanço das práticas de projeto de fundações, promovendo maior eficiência, sustentabilidade e racionalização no uso dos recursos estruturais.